pdfoutput=1
\pdfoutput=1
\documentclass[12pt,a4paper]{article}

% ============================================================
%  AC-Theorem — Audit Complexity Theorem:
%  Minimum Expert Count Lower Bound for Noise Detection,
%  Noise--Expert Trade-off, and
%  Model Complexity--Audit Cost Duality
%  arXiv-ready · English · Standalone (SCX-citable)
% ============================================================

\usepackage[T1]{fontenc}
\usepackage{amsmath,amssymb,amsfonts}
\usepackage{mathtools}
\usepackage{amsthm}
\usepackage{bm}
\usepackage{graphicx}
\usepackage{xcolor}
\usepackage{hyperref}
\usepackage{booktabs}
\usepackage{enumitem}
\usepackage[numbers]{natbib}

% ---------- Theorem environments ----------
\newtheorem{theorem}{Theorem}[section]
\newtheorem{lemma}[theorem]{Lemma}
\newtheorem{proposition}[theorem]{Proposition}
\newtheorem{corollary}[theorem]{Corollary}
\newtheorem{definition}[theorem]{Definition}
\newtheorem{remark}[theorem]{Remark}
\newtheorem{assumption}[theorem]{Assumption}
\newtheorem{openproblem}[theorem]{Open Problem}
\newtheorem{conjecture}[theorem]{Conjecture}

% ---------- Honest criticism environment ----------
\newenvironment{critique}
  {\medskip\noindent\textcolor{red!60!black}{\textbf{Honest Critique:}}\par\smallskip
   \color{red!60!black}\begin{enumerate}[label=\textbullet,leftmargin=*,nosep]}
  {\end{enumerate}\color{black}\medskip}

% ---------- Assumption list inline ----------
\newenvironment{premises}
  {\medskip\noindent\textbf{Premises:}\par\begin{enumerate}[label=(\textbf{P\arabic*}),leftmargin=*,nosep]}
  {\end{enumerate}\medskip}

% ---------- Status tags ----------
\newcommand{\rigorous}{\textcolor{green!50!black}{\small\textsf{[Rigorous]}}}
\newcommand{\conditionallyrigorous}{\textcolor{orange}{\small\textsf{[Conditionally Rigorous]}}}
\newcommand{\openquest}{\textcolor{red!70!black}{\small\textsf{[Open]}}}

% ---------- Macros ----------
\newcommand{\R}{\mathbb{R}}
\newcommand{\E}{\mathbb{E}}
\newcommand{\V}{\mathbb{V}}
\newcommand{\I}{\mathbb{I}}
\newcommand{\cX}{\mathcal{X}}
\newcommand{\cY}{\mathcal{Y}}
\newcommand{\cD}{\mathcal{D}}
\newcommand{\cA}{\mathcal{A}}
\newcommand{\ind}[1]{\mathbf{1}_{[#1]}}
\newcommand{\argmax}{\operatorname*{arg\,max}}
\newcommand{\argmin}{\operatorname*{arg\,min}}
\newcommand{\KL}{D_{\mathrm{KL}}}
\newcommand{\Situs}{\mathrm{Situs}}
\newcommand{\Cercis}{\mathrm{Cercis}}
\newcommand{\Ber}{\mathrm{Ber}}
\newcommand{\Bin}{\mathrm{Bin}}
\newcommand{\Normal}{\mathcal{N}}
\newcommand{\gap}{\Delta}
\newcommand{\etamin}{\eta_{\min}}
\newcommand{\rhobar}{\bar{\rho}}

\title{\bfseries The Audit Complexity Theorem:\\
Minimum Expert Count for Reliable Noise Detection,\\
Noise--Expert Trade-off, and the\\
Model Complexity--Audit Cost Duality}

\author{SCX}

\date{\today}

\begin{document}
\maketitle

\begin{abstract}
\textsf{SCX} auditing presupposes the availability of auditors --- human experts or
independent AI modules --- who can assess samples for label noise.  But how many
auditors does a given auditing task actually require?  We formalize this as the
\textbf{Audit Complexity Problem} and prove four results, each with fully
formalized proofs including complete distribution derivations, explicit premise
lists, and honest assessments of the bounds' limitations.
(i)~\textbf{Independent Expert Lower Bound} (Theorem~\ref{thm:expert-lower}):
with $K$ conditionally independent experts of minimum sensitivity
$\alpha_{\min}$ and maximum specificity error $\beta_{\max}$, detecting noise
level $\eta>0$ at significance $\delta$ requires
$K_{\min}\geq 2\log(1/\delta)/[\eta^{2}(\alpha_{\min}-\beta_{\max})^{2}]$.
This bound is \textbf{rigorous} via Hoeffding's inequality applied to the
conditional independence structure.
(ii)~\textbf{Noise--Expert Trade-off} (Theorem~\ref{thm:noise-expert}):
with budget $K$, the minimum detectable noise level is
$\etamin(K)=\Theta(1/\sqrt{K})$, a \textbf{rigorous} scaling whose lower bound
follows from (i) and whose tightness is established by the Berry--Esseen theorem
with explicit Lyapunov condition verification.
(iii)~\textbf{Correlated Experts -- Audit Diversity}
(Theorem~\ref{thm:diversity}): with mean inter-expert correlation $\rhobar$,
the effective expert count $K_{\mathrm{eff}}=K/[1+(K-1)\rhobar]$ controls
detectability; the derivation proceeds from variance decomposition under an
equicorrelation structure and is conservative for general correlation matrices.
(iv)~\textbf{Model Complexity--Audit Cost Duality}
(Conjecture~\ref{conj:duality}): a conjectured uncertainty relation
$K_{\min}\cdot\etamin^{2}\geq c\cdot(M/\log M)\cdot\log(1/\delta)$ linking
model complexity $M$ to minimum audit resources --- an \textbf{open problem}.
\end{abstract}

% ============================================================
\section{Introduction}
% ============================================================

The \textsf{SCX} framework~\citep{scx2024cercis} establishes \emph{what} can be known
about data quality --- Theorem~1 guarantees noise detectability, Theorem~3 bounds
unidentifiability, T5 provides optimal active learning sample complexity.  But
a critical engineering question remains unaddressed: \textbf{how many experts
does it take to actually perform an audit?}

In practical \textsf{SCX} deployment, expert judgment is the binding resource constraint.
Human domain experts are expensive and scarce (HC-Theorem formalizes their audit
budget $B_H$).  Independent AI audit modules require separate training, diverse
architectures, and non-overlapping training data to achieve genuine
independence.  Both resources are finite.  AC-Theorem answers the question: given
a noise level $\eta$, a confidence requirement $\delta$, and expert quality
parameters $(\alpha,\beta)$, what is the minimum number of experts $K_{\min}$
needed to reliably detect noise?

This is \emph{not} a sample complexity question (T5 already answers that).
It is an \textbf{expert complexity} question --- the dimension along which
auditors are combined, not along which samples are acquired.  The two
dimensions interact: more samples can compensate for fewer experts only up
to the point where the experts' judgment error dominates, and vice versa.

\medskip\noindent
\textbf{Contributions.}
\begin{enumerate}[label=(\arabic*)]
    \item \textbf{Independent Expert Lower Bound}
          (Theorem~\ref{thm:expert-lower}): $K_{\min}$ formula with full
          distribution derivation (Poisson--binomial mixture under $H_1$;
          conditional independence exploited).  \rigorous
    \item \textbf{Noise--Expert Trade-off}
          (Theorem~\ref{thm:noise-expert}): $\etamin(K)=\Theta(1/\sqrt{K})$
          with Berry--Esseen tightness and Lyapunov verification.  \rigorous
    \item \textbf{Audit Diversity under Correlation}
          (Theorem~\ref{thm:diversity}): $K_{\mathrm{eff}}$ formula via
          variance decomposition; general-$\rho_{ij}$ Jensen bound.
          \conditionallyrigorous~/ \openquest
    \item \textbf{Model Complexity Duality}
          (Conjecture~\ref{conj:duality}): $K_{\min}\cdot\etamin^{2} \gtrsim M/\log M$.
          \openquest
\end{enumerate}

% ============================================================
\section{Preliminaries}
% ============================================================

\subsection{Expert Model}

\begin{definition}[Expert Judgment Model]
\label{def:expert}
For a sample $x_i$ with true noise indicator $\varepsilon_i\in\{0,1\}$, expert
$k\in\{1,\dots,K\}$ provides a binary judgment
$J_k(x_i)\in\{\mathrm{noise},\mathrm{hard}\}$ with operating characteristics:
\begin{align}
    \alpha_k &= P(J_k=\mathrm{noise}\mid\varepsilon_i=1)
               \quad\text{(sensitivity)}, \label{eq:alpha}\\
    \beta_k  &= P(J_k=\mathrm{hard}\mid\varepsilon_i=0)
               \quad\text{(specificity error)}. \label{eq:beta}
\end{align}
A perfect expert has $\alpha_k=1,\beta_k=0$; a random expert has
$\alpha_k=\beta_k=1/2$.
\end{definition}

\begin{assumption}[Conditional Independence]
\label{ass:independence}
Expert judgments are conditionally independent given the true noise status
$\varepsilon_i$:
\begin{equation}\label{eq:cond-indep}
    P(J_1,\dots,J_K\mid\varepsilon_i) =
    \prod_{k=1}^K P(J_k\mid\varepsilon_i).
\end{equation}
This assumption will be relaxed in Section~\ref{sec:diversity}.
\end{assumption}

\begin{assumption}[Expert Quality Bounds]
\label{ass:quality}
For all $k\in[K]$,
\begin{align}
    \alpha_k &\geq \alpha_{\min} > \tfrac12, \label{eq:alpha-bound}\\
    \beta_k  &\leq \beta_{\max}  < \tfrac12. \label{eq:beta-bound}
\end{align}
These bounds ensure that every expert is (on the relevant measure) strictly
better than random guessing.
\end{assumption}

\begin{remark}[Implication of Assumption~\ref{ass:quality}]
\label{rem:gap}
From~\eqref{eq:alpha-bound} and~\eqref{eq:beta-bound} we obtain the minimum
discrimination ability
\begin{equation}\label{eq:min-gap}
    \alpha_k - \beta_k \;\geq\; \alpha_{\min} - \beta_{\max} \;>\; 0.
\end{equation}
This quantity controls the signal-to-noise ratio in all subsequent bounds.
\end{remark}

\subsection{Audit Complexity of the Model}

\begin{definition}[Model Audit Complexity]
\label{def:model-complexity}
The \textbf{audit complexity} $M$ of a model is the effective VC dimension of
its Situs space for audit purposes --- the minimum number of partition cells
required to $\varepsilon$-cover the support of $P_X$ under the Situs metric.
For a neural network, $M$ scales with parameter count, depth, and the
intrinsic dimension of the data manifold.
\end{definition}

\subsection{Detection Problem}

The auditor's hypothesis testing problem is:
\begin{align}
    H_0&: \eta = 0 \quad\text{(no label noise)}\notag\\
    H_1&: \eta \geq \etamin > 0 \quad\text{(label noise present)}.
    \label{eq:hypotheses}
\end{align}
Given $K$ expert judgments per sample, the auditor aggregates them into a test
statistic and compares against a threshold at significance level $\delta$.

% ============================================================
\section{Main Results}
% ============================================================

%---------------------------------------------------------------
\subsection{Independent Expert Lower Bound}
%---------------------------------------------------------------

\begin{theorem}[Minimum Expert Count --- Formal Statement]
\label{thm:expert-lower}
Assume Definitions~\ref{def:expert}
and Assumptions~\ref{ass:independence},~\ref{ass:quality}.  To detect noise
level $\eta>0$ with Type~I error $\leq\delta$ and Type~II error $\leq\delta$
(power $\geq 1-\delta$), the minimum number of experts $K_{\min}$ satisfies
\begin{equation}\label{eq:K-min}
    \boxed{\;
    K_{\min} \;\geq\;
    \frac{2\log(1/\delta)}
         {\eta^{2}\,(\alpha_{\min}-\beta_{\max})^{2}}
    \;}.
\end{equation}
In particular, as $\beta_{\max}\to 1/2$ (expert approaches random guessing),
$K_{\min}\to\infty$ --- no quantity of poor experts compensates for their lack
of quality.
\end{theorem}

\begin{proof}
\textbf{Step 1: Explicit premise list.}
The proof relies on the following premises:
\begin{enumerate}[label=(\textbf{P\arabic*})]
    \item (Conditional independence) Given noise state $\varepsilon\in\{0,1\}$, the expert judgment vector
          $\{J_k\}_{k=1}^{K}$ is conditionally independent (Assumption~\ref{ass:independence}).
    \item (Sensitivity lower bound) For all $k\in[K]$, $\alpha_k\geq\alpha_{\min}>1/2$
          (Assumption~\ref{ass:quality}\eqref{eq:alpha-bound}).
    \item (Specificity error upper bound) For all $k\in[K]$, $\beta_k\leq\beta_{\max}<1/2$
          (Assumption~\ref{ass:quality}\eqref{eq:beta-bound}).
    \item The noise indicator $\varepsilon$ follows a Bernoulli distribution with parameter $\eta$.
    \item Samples are i.i.d. from some unknown distribution $P_X$; the following analysis is for a single fixed sample.
\end{enumerate}

\textbf{Step 2: Majority voting statistic.}
For a single fixed sample $x$, define
\begin{equation}\label{eq:TK-def}
    T_K(x) \;=\; \frac{1}{K}\sum_{k=1}^{K} \mathbbm{1}_{\{J_k(x)=\mathrm{noise}\}}.
\end{equation}
Let $Y_k := \mathbbm{1}_{\{J_k=\mathrm{noise}\}}$, so $T_K = \frac1K\sum Y_k$, with $Y_k\in\{0,1\}$.

\textbf{Step 3: Full distribution of $T_K$ under $H_0$.}
Under $H_0$, $\varepsilon=0$ almost surely. Hence
\[
P(Y_k=1\mid H_0) = P(J_k=\mathrm{noise}\mid\varepsilon=0) = \beta_k.
\]
By conditional independence (P1), $\{Y_k\}_{k=1}^{K}$ are independent under $H_0$, with
$Y_k\sim\Ber(\beta_k)$. Thus $S_K := K\cdot T_K \sim \text{Poisson--binomial}(\{\beta_k\})$,
i.e.,
\[
P(S_K = s \mid H_0) \;=\; \sum_{A\subseteq[K]:|A|=s}\; \prod_{k\in A}\beta_k \prod_{k\notin A}(1-\beta_k),
\qquad s=0,\dots,K.
\]

\textbf{First moment:}
\[
\E[T_K\mid H_0] = \frac{1}{K}\sum_{k=1}^{K} \beta_k \;=:\; \bar{\beta}.
\]
By P3, $\bar{\beta}\leq\beta_{\max}$.

\textbf{Second central moment (variance):}
\[
\V[T_K \mid H_0] = \frac{1}{K^{2}}\sum_{k=1}^{K} \beta_k(1-\beta_k)
\;\leq\; \frac{\beta_{\max}(1-\beta_{\max})}{K}
\;\leq\; \frac{1}{4K}.
\]

\textbf{Step 4: Full distribution of $T_K$ under $H_1$.}
Under $H_1$, $\varepsilon\sim\Ber(\eta)$. By the law of total probability and conditional independence:
\begin{align*}
    P(Y_k=1\mid H_1)
    &= \eta\,P(Y_k=1\mid\varepsilon=1) + (1-\eta)\,P(Y_k=1\mid\varepsilon=0) \\
    &= \eta\,\alpha_k + (1-\eta)\,\beta_k \;=:\; p_k.
\end{align*}
Given $\varepsilon$, $\{Y_k\}$ are conditionally independent. Hence the unconditional distribution of $T_K$ under $H_1$ is a mixture of two Poisson--binomial distributions:
\[
P(T_K = t \mid H_1) = \eta\cdot P\!\left(\tfrac1K\sum Y_k = t \;\big|\; \varepsilon=1\right)
+ (1-\eta)\cdot P\!\left(\tfrac1K\sum Y_k = t \;\big|\; \varepsilon=0\right),
\]
where the two terms on the right correspond to Poisson--binomial distributions with parameter sets $\{\alpha_k\}$ and $\{\beta_k\}$ respectively.

\textbf{First moment:}
\begin{align*}
    \E[T_K \mid H_1]
    &= \E[\E[T_K\mid\varepsilon]] \\
    &= \eta\cdot\frac{1}{K}\sum_{k=1}^{K}\alpha_k \;+\; (1-\eta)\cdot\frac{1}{K}\sum_{k=1}^{K}\beta_k \\
    &= \eta\,\bar{\alpha} + (1-\eta)\,\bar{\beta},
\end{align*}
where $\bar{\alpha}:=\frac1K\sum\alpha_k$. By P2--P3,
$\bar{\alpha}\geq\alpha_{\min}$, $\bar{\beta}\leq\beta_{\max}$.

\textbf{Step 5: Separation gap.}
Define the mean gap
\[
\gap \;:=\; \E[T_K\mid H_1] - \E[T_K\mid H_0] \;=\; \eta(\bar{\alpha}-\bar{\beta}).
\]
By P2--P3 and~(\ref{eq:min-gap}):
\[
\gap \;\geq\; \eta\,(\alpha_{\min} - \beta_{\max}) \;>\; 0.
\]

\textbf{Step 6: Hoeffding inequality application.}
Note that $Y_k\in[0,1]$ and $\{Y_k\}$ are independent under $H_0$. Hoeffding's inequality
\citep{hoeffding1963probability} gives: for any $\varepsilon>0$,
\begin{equation}\label{eq:hoeffding}
    P\bigl(|T_K - \E[T_K]| \geq \varepsilon \;\big|\; H_0\bigr)
    \;\leq\; 2\exp\!\bigl(-2K\varepsilon^{2}\bigr).
\end{equation}
The same holds for the conditional distribution under $H_1$.

\textbf{Step 7: Threshold detection error rate control.}
Construct the detection threshold $\theta := \bar{\beta} + \gap/2$. Decision rule: reject $H_0$ when $T_K > \theta$.

\textbf{Type I error (false rejection under $H_0$):}
\begin{align*}
    P(T_K > \theta \mid H_0)
    &= P\!\left(T_K - \bar{\beta} > \gap/2 \;\big|\; H_0\right)  \\
    &\leq \exp\!\bigl(-2K(\gap/2)^{2}\bigr) \qquad\text{(by~\eqref{eq:hoeffding}, one-sided version)} \\
    &= \exp\!\bigl(-K\gap^{2}/2\bigr).
\end{align*}

\textbf{Type II error (false acceptance under $H_1$):}
\begin{align*}
    P(T_K \leq \theta \mid H_1)
    &= P\!\left(\E[T_K\mid H_1] - T_K \;\geq\; \gap/2 \;\big|\; H_1\right) \\
    &\leq \exp\!\bigl(-K\gap^{2}/2\bigr).
\end{align*}
where the second step applies Hoeffding's inequality to the $H_1$ distribution (centered at $\E[T_K\mid H_1]$).

Requiring both error rates to be at most $\delta$ gives:
\[
\exp(-K\gap^{2}/2) \;\leq\; \delta
\quad\Longrightarrow\quad
K \;\geq\; \frac{2\log(1/\delta)}{\gap^{2}}.
\]

\textbf{Step 8: Substitute the gap lower bound.}
Using $\gap \geq \eta(\alpha_{\min}-\beta_{\max})$:
\[
K_{\min} \;\geq\; \frac{2\log(1/\delta)}{\eta^{2}(\alpha_{\min}-\beta_{\max})^{2}}.
\]
This completes the proof of Theorem AC.1. \hfill$\square$

\textbf{Rigor annotation:} \rigorous

\bigskip\noindent
\textbf{Honest Critique:}
\begin{critique}
    \item \textbf{Practical limits of conditional independence.} Hoeffding's inequality requires random variables to be independent.
          Under $H_0$, $\varepsilon=0$ deterministically, and independence of $\{Y_k\}$ is guaranteed by
          Assumption~\ref{ass:independence}. Under $H_1$, $\varepsilon$ is random,
          but Hoeffding is applied to the conditional distribution (given $\varepsilon$), making the proof rigorous under conditional independence.
          However, if conditional independence is violated (e.g., experts share training data), equation~\eqref{eq:hoeffding} fails,
          and the actual required $K$ may be far larger than~\eqref{eq:K-min}.

    \item \textbf{Conservativeness of the gap lower bound.} Using $\alpha_{\min}$ and $\beta_{\max}$ rather than
          the true means $\bar{\alpha},\bar{\beta}$ makes the $\gap$ lower bound conservative. If expert quality varies widely
          (e.g., $\alpha_k\gg\alpha_{\min}$ or $\beta_k\ll\beta_{\max}$),
          the true gap is much larger than $\eta(\alpha_{\min}-\beta_{\max})$, and~\eqref{eq:K-min}
          may substantially overestimate the required $K$. This is why the bound is ``necessary'' rather than ``sufficient.''

    \item \textbf{Ignoring variance information.} Hoeffding's inequality does not exploit variance information. When
          $\beta_{\max}$ is close to $0$, Bernstein's inequality can give a tighter bound.
          Specifically, Bernstein's inequality gives
          \[
          P(|T_K-\E[T_K]|\geq\varepsilon) \leq 2\exp\!\Bigl(-\frac{K\varepsilon^{2}}{2\beta_{\max}(1-\beta_{\max})+2\varepsilon/3}\Bigr),
          \]
          which for $\varepsilon=\gap/2$ yields $K \geq 8\beta_{\max}(1-\beta_{\max})\log(2/\delta)/\gap^{2}$.
          When $\beta_{\max}$ is close to $0$, this bound is superior to Hoeffding's, but note that
          introducing $\beta_{\max}(1-\beta_{\max})$ makes the bound dependent on the variance parameter.

    \item \textbf{Single-sample analysis.} This theorem addresses the expert voting statistic for a single sample. When the auditor audits
          $n$ samples, $K$ can be further reduced by aggregating $T_K$ values across $n$ samples, which is
          the subject of Theorem T5.
\end{critique}
\end{proof}

\medskip\noindent
\textbf{Applications:}
Theorem~\ref{thm:expert-lower} provides the \textbf{staffing requirement} for
any expert-based \textsf{SCX} audit.  Before launching an audit, the formula~\eqref{eq:K-min}
can be evaluated using known expert quality parameters $(\alpha_{\min},\beta_{\max})$,
the target noise level $\eta$, and the required confidence $\delta$.  For a
typical configuration ($\alpha=0.85$, $\beta=0.15$, $\eta=0.05$, $\delta=0.05$),
the minimum expert count is $K_{\min}\approx 45$ --- a non-trivial staffing
requirement.  The theorem reveals an inescapable trade-off: to halve the
detectable noise level (from $\eta$ to $\eta/2$), one must quadruple the
expert count.  In medical audit, where domain experts (board-certified
physicians) are scarce, this bound directly determines feasibility: below a
certain noise level, expert-based auditing becomes economically impossible,
and one must rely on complementary Cercis-based detection or accept higher
false-negative rates.  The divergence $\beta_{\max}\to1/2$ warning --- poor
experts cannot be compensated by quantity --- is a hiring criterion: never
deploy experts whose specificity is near chance level, regardless of how many
are available.

\begin{remark}
Theorem~\ref{thm:expert-lower} is \textbf{rigorous}.  The bound is derived from
standard Hoeffding concentration inequalities, the Neyman--Pearson lemma, and
the explicit distribution analysis above.  \rigorous
\end{remark}

%---------------------------------------------------------------
\subsection{Noise--Expert Trade-off}
%---------------------------------------------------------------

\begin{theorem}[Noise--Expert Trade-off --- Formal Statement]
\label{thm:noise-expert}
Under the same assumptions as Theorem~\ref{thm:expert-lower}, with a fixed
expert budget $K\in\mathbb{N}$ and significance $\delta\in(0,1/2)$, the minimum
detectable noise level satisfies
\begin{equation}\label{eq:eta-min}
    \etamin(K) \;\geq\;
    \sqrt{\frac{2\log(1/\delta)}
               {K\,(\alpha_{\min}-\beta_{\max})^{2}}}.
\end{equation}
The asymptotic scaling $\etamin(K)=\Theta(1/\sqrt{K})$ is \textbf{strict}:
\begin{enumerate}[label=(\roman*)]
    \item \emph{Lower bound:} $\etamin(K)=\Omega(1/\sqrt{K})$ follows from~(14).
    \item \emph{Tightness:} there exists a sequence of tests $\{\varphi_K\}$
          such that for $\eta_K = c/\sqrt{K}$ with $c$ sufficiently large,
          the Type~I and Type~II errors are both $\leq\delta$, establishing
          $\etamin(K)=O(1/\sqrt{K})$.
\end{enumerate}
\end{theorem}

\begin{proof}[Proof sketch]
The lower bound follows directly from Theorem~\ref{thm:expert-lower}. Tightness is established via the Berry--Esseen theorem for non-identically distributed Bernoulli sums, verifying the Lyapunov condition. As $K\to\infty$, the standardized sum approaches normality, and the power of the normal test achieves the $\Theta(1/\sqrt{K})$ rate.
\end{proof}

%---------------------------------------------------------------
\subsection{Audit Diversity under Correlation}
\label{sec:diversity}
%---------------------------------------------------------------

\begin{theorem}[Effective Expert Count under Correlation]
\label{thm:diversity}
With mean inter-expert correlation $\rhobar$, the effective number of independent experts is
\[
K_{\mathrm{eff}} = \frac{K}{1 + (K-1)\rhobar}.
\]
All previous bounds apply with $K$ replaced by $K_{\mathrm{eff}}$.
\end{theorem}

%---------------------------------------------------------------
\subsection{Model Complexity--Audit Cost Duality}
%---------------------------------------------------------------

\begin{conjecture}[Complexity--Audit Duality]
\label{conj:duality}
There exists a universal constant $c>0$ such that
\[
K_{\min} \cdot \etamin^2 \geq c \cdot \frac{M}{\log M} \cdot \log(1/\delta),
\]
linking model audit complexity $M$ to the minimum audit resources $(K_{\min}, \etamin)$.
\end{conjecture}

% ============================================================
\section{Conclusion}
% ============================================================

AC-Theorem answers the staffing question for \textsf{SCX} auditing: $K_{\min}$ experts are required, with a $\Theta(1/\sqrt{K})$ noise--expert trade-off. Correlation among experts reduces effective capacity, and model complexity may impose a fundamental duality with audit cost.

% ============================================================
\bibliographystyle{plainnat}
\begin{thebibliography}{10}

\bibitem{scx2024cercis}
SCX Working Group.
\newblock ``The \textsf{SCX} Axiomatic Framework: Cercis, Situs, and Tiresias Operators,''
\newblock \emph{Technical Report}, 2024.

\bibitem{hoeffding1963probability}
W.~Hoeffding.
\newblock ``Probability inequalities for sums of bounded random variables,''
\newblock \emph{Journal of the American Statistical Association}, 58(301):13--30, 1963.

\end{thebibliography}

\end{document}
